\documentclass[11pt]{article}
\usepackage{mathblatt}
\usepackage{adjustbox,varwidth}
% gesetzt von werkzeuge/setzer.py aus dem Lernweg (L4-A · L4-B · L4-C · L4-D · L4-T) und den Bankzeilen; Muster T6B.
\input{praeambel}
\newcommand{\kennung}{4CR}
\newcommand{\grau}[1]{{\color{mbgrau}#1}}

\begin{document}
\pfheftstil
\fancyfoot[C]{\parbox[t]{\textwidth}{\mbfhbox\par\vspace{1.5mm}{\footnotesize\color{mbgrau}{\scriptsize \kennung}\hfill\thepage}}}
\pfheftkopf{Kenngrößen}{Klasse 7}
% L4-A: Spannweite und Modalwert
\begin{pfaufg}{}{1.}{}
\grau{a)\ Ben springt siebenmal weit.\par\vspace{3pt} \adjustbox{max width=\linewidth}{\begin{varwidth}{50cm}\sachtabelle{lccccccc}{Versuch & $1$ & $2$ & $3$ & $4$ & $5$ & $6$ & $7$}{Weite in m & $3{,}8$ & $4{,}1$ & $3{,}6$ & $4{,}1$ & $3{,}9$ & $4{,}1$ & $3{,}7$}\end{varwidth}}\par\vspace{3pt} Bestimme Minimum, Maximum, Spannweite und Modalwert. Tom sagt: „Die Spannweite ist $3{,}6$ m bis $4{,}1$ m.“ Stimmt das?\par\smallskip Ordnen: $3{,}6$; $3{,}7$; $3{,}8$; $3{,}9$; $4{,}1$; $4{,}1$; $4{,}1$.\par Minimum: vorne, $3{,}6$\,m. Maximum: hinten, $4{,}1$\,m.\par Spannweite $= 4{,}1\,\text{m} - 3{,}6\,\text{m} = 0{,}5$\,m.\par Modalwert: $4{,}1$\,m kommt dreimal vor, öfter als jeder andere Wert.\par Tom prüfen: Die Spannweite ist ein Abstand, also eine Zahl. Toms Aussage ist falsch; richtig ist $0{,}5$\,m.}
\par\medskip
b)\ Mia schreibt eine Woche lang auf, wie viele Minuten sie für die Hausaufgaben braucht: $25$, $40$, $15$, $30$, $45$, $20$, $35$. Bestimme Minimum, Maximum und Spannweite.\karo{3}
\fusshilfe{\mbox{1:~b) Min.\ $15$, Max.\ $45$, Spannweite $30$\,min}}
\end{pfaufg}

% L4-A: Spannweite und Modalwert
\begin{pfaufg}{}{2.}{}
Sechs Bäckereien verlangen für ein Brötchen (in €): $0{,}45$; $0{,}50$; $0{,}40$; $0{,}55$; $0{,}50$; $0{,}60$. Bestimme die Spannweite und den Modalwert.\karo{3}
\fusshilfe{\mbox{2:~Spannweite $0{,}20$\,€, Modalwert $0{,}50$\,€}}
\end{pfaufg}

% L4-A: Spannweite und Modalwert
\begin{pfaufg}{}{3.}{}
Die Tabelle zeigt die Preise einer Tankstelle in € je Liter.\par\vspace{3pt} \adjustbox{max width=\linewidth}{\begin{varwidth}{50cm}\sachtabelle{lccccc}{ & Mo & Di & Mi & Do & Fr}{Super & $1{,}79$ & $1{,}84$ & $1{,}76$ & $1{,}81$ & $1{,}85$\\ Diesel & $1{,}69$ & $1{,}66$ & $1{,}72$ & $1{,}70$ & $1{,}68$}\end{varwidth}}\par\vspace{3pt} Bei welcher Sorte schwankte der Preis in dieser Woche stärker?\karo{3}
\fusshilfe{\mbox{3:~Super ($0{,}09$\,€ gegen $0{,}06$\,€)}}
\end{pfaufg}

% L4-A: Spannweite und Modalwert
\begin{pfaufg}{}{4.}{}
Die Kinder einer Schwimmgruppe sind so alt (in Jahren): $9$, $11$, $10$, $12$, $9$, $13$, $10$, $9$, $11$, $12$.\par Aussage 1: „Die Spannweite beträgt $4$ Jahre.“\par Aussage 2: „Der Modalwert beträgt $13$ Jahre.“\par Eine Aussage ist falsch. Kreuze sie an und berichtige sie.\par\smallskip \gkreuz{Aussage 1}\gkreuz{Aussage 2}
\fusshilfe{\mbox{4:~Aussage 2; richtig: Modalwert $9$ Jahre}}
\end{pfaufg}

% L4-B: Median
\begin{pfaufg}{}{5.}{}
\grau{a)\ Bestimme den Median. a)~Lena braucht an sieben Tagen so viele Minuten zur Schule: $12$, $8$, $15$, $10$, $8$, $20$, $11$. b)~Ole braucht an sechs Tagen: $18$, $21$, $16$, $19$, $23$, $17$.\par\smallskip a) Ordnen: $8,\ 8,\ 10,\ 11,\ 12,\ 15,\ 20$.\par Von beiden Enden je einen Wert streichen, bis einer übrig bleibt: Bei $7$ Werten bleibt der $4.$ Wert. Median $= 11$\,min.\par b) Ordnen: $16,\ 17,\ 18,\ 19,\ 21,\ 23$.\par Bei $6$ Werten bleiben zwei in der Mitte übrig: $18$ und $19$.\par Median $= (18 + 19) : 2 = 37 : 2 = 18{,}5$\,min.}
\par\medskip
b)\ Fünf Kinder werfen auf einen Korb. Sie treffen so oft: $7$, $3$, $9$, $4$, $6$. Bestimme den Median.\karo{3}
\fusshilfe{\mbox{5:~b) $6$ Treffer}}
\end{pfaufg}

% L4-B: Median
\begin{pfaufg}{}{6.}{}
Sechs Pflanzen sind so hoch (in cm): $14$, $9$, $12$, $20$, $11$, $16$. Bestimme den Median.\karo{3}
\fusshilfe{\mbox{6:~$13$\,cm}}
\end{pfaufg}

% L4-B: Median
\begin{pfaufg}{}{7.}{}
Die Tabelle zeigt die höchste Temperatur an sechs Tagen.\par\vspace{3pt} \adjustbox{max width=\linewidth}{\begin{varwidth}{50cm}\sachtabelle{lcccccc}{Tag & Mo & Di & Mi & Do & Fr & Sa}{in °C & $23$ & $17$ & $20$ & $25$ & $19$ & $18$}\end{varwidth}}\par\vspace{3pt} Bestimme den Median.\karo{3}
\fusshilfe{\mbox{7:~$19{,}5$\,°C}}
\end{pfaufg}

% L4-B: Median
\begin{pfaufg}{}{8.}{}
Tim hat in sieben Basketballspielen $12$, $18$, $9$, $15$, $21$, $14$, $10$ Punkte gemacht, Ali in sechs Spielen $16$, $8$, $13$, $22$, $11$, $17$ Punkte. Der Trainer vergleicht die Mediane der beiden. Wer liegt vorn?\karo{3}
\fusshilfe{\mbox{8:~Ali ($14{,}5$ gegen $14$)}}
\end{pfaufg}

% L4-C: Arithmetisches Mittel
\begin{pfaufg}{}{9.}{}
\gzwei{\adjustbox{max width=\linewidth}{\begin{varwidth}{50cm}\saeulenab[ymax=80,ystep=20,yfein=10,ylabel=Kunden,hoehe=3.2,breite=0.8]{Mo/40,Di/50,Mi/30,Do/60,Fr/50,Sa/70,So/60}\end{varwidth}}}{\grau{a)\ Das Säulendiagramm zeigt, wie viele Kunden an jedem Tag einer Woche an einem Kiosk eingekauft haben. Wie viele Kunden kamen im Durchschnitt pro Tag?\par\smallskip Ablesen: Von $0$ bis $20$ sind es $2$ Kästchen, ein Kästchen ist $10$ Kunden.\par Werte: Mo $40$, Di $50$, Mi $30$, Do $60$, Fr $50$, Sa $70$, So $60$.\par Anzahl: $7$ Tage, also $7$ Werte.\par Summe: $40 + 50 + 30 + 60 + 50 + 70 + 60 = 360$.\par Teilen: $360 : 7 = 51{,}42\ldots$\par Runden: Geht die Division nicht auf, runde: bei Personen auf eine ganze Zahl, sonst auf eine Stelle nach dem Komma. Kunden zählt man ganz, im Durchschnitt kamen etwa $51$ Kunden pro Tag.}}
\par\medskip
b)\ Ein Team schießt in fünf Spielen $3$, $1$, $4$, $0$, $2$ Tore. Berechne den Mittelwert.\karo{3}
\fusshilfe{\mbox{9:~b) $2$ Tore pro Spiel}}
\end{pfaufg}

% L4-C: Arithmetisches Mittel
\begin{pfaufg}{}{10.}{}
Pia springt viermal weit: $4{,}12$ m; $3{,}95$ m; $4{,}30$ m; $4{,}07$ m. Berechne ihre durchschnittliche Weite.\karo{3}
\fusshilfe{\mbox{10:~$4{,}11$\,m}}
\end{pfaufg}

% L4-C: Arithmetisches Mittel
\begin{pfaufg}{}{11.}{}
Zu den $17$ Heimspielen eines Vereins kamen in einer Saison zusammen $412\,390$ Zuschauer. Wie viele Zuschauer kamen im Durchschnitt zu einem Spiel?\karo{3}
\fusshilfe{\mbox{11:~$\approx 24\,258$ Zuschauer}}
\end{pfaufg}

% L4-C: Arithmetisches Mittel
\begin{pfaufg}{}{12.}{}
\gzwei{\adjustbox{max width=\linewidth}{\begin{varwidth}{50cm}\saeulenab[ymax=350,ystep=100,yfein=50,ylabel=Besucher,hoehe=3.2,breite=0.8]{Mo/150,Di/100,Mi/150,Do/150,Fr/200,Sa/300,So/300}\end{varwidth}}}{%
Das Säulendiagramm zeigt die Besucher eines Freibads in einer Woche. Der Bademeister sagt: „Im Durchschnitt kamen mehr als $200$ Besucher am Tag.“ Hat er recht? Begründe mit einer Rechnung.\karo{3}}
\fusshilfe{\mbox{12:~Nein, nur $\approx 193$ Besucher am Tag}}
\end{pfaufg}

% L4-D: Mittel rückwärts und Änderungen
\begin{pfaufg}{}{13.}{}
\grau{a)\ Die Tabelle zeigt die Mittagstemperatur einer Woche. Der Sonntag fehlt. Im Mittel waren es $20$ °C.\par\vspace{3pt} \adjustbox{max width=\linewidth}{\begin{varwidth}{50cm}\sachtabelle{lccccccc}{Tag & Mo & Di & Mi & Do & Fr & Sa & So}{in °C & $18$ & $21$ & $19$ & $23$ & $20$ & $22$ & ?}\end{varwidth}}\par\vspace{3pt} a)~Wie warm war es am Sonntag? b)~Lisa streicht den wärmsten und den kältesten Tag. Wie ändert sich das Mittel?\par\smallskip a) Summe aller $7$ Werte $=$ Mittel $\cdot$ Anzahl $= 20 \cdot 7 = 140$.\par Bekannte Werte: $18 + 21 + 19 + 23 + 20 + 22 = 123$.\par Sonntag $= 140 - 123 = 17$\,°C. Probe: $(123 + 17) : 7 = 20$.\par b) Wärmster Tag $23$\,°C, kältester $17$\,°C. Neue Summe: $140 - 23 - 17 = 100$. Neue Anzahl: $7 - 2 = 5$.\par Neues Mittel: $100 : 5 = 20$\,°C. Es bleibt gleich, weil die zwei gestrichenen Werte zusammen $40 = 2 \cdot 20$ sind.}
\par\medskip
b)\ Drei Zahlen haben den Mittelwert $8$. Zwei davon sind $5$ und $9$. Wie heißt die dritte Zahl?\karo{3}
\fusshilfe{\mbox{13:~b) $10$}}
\end{pfaufg}

% L4-D: Mittel rückwärts und Änderungen
\begin{pfaufg}{}{14.}{}
Tom läuft vier Runden. Im Mittel braucht er $62$ s für eine Runde. Die ersten drei Runden dauern $60$ s, $65$ s und $59$ s. Wie lange dauert die vierte Runde?\karo{3}
\fusshilfe{\mbox{14:~$64$\,s}}
\end{pfaufg}

% L4-D: Mittel rückwärts und Änderungen
\begin{pfaufg}{}{15.}{}
Fünf Kinder einer Gruppe sind im Mittel $12$ Jahre alt. Das älteste Kind ist $16$ Jahre alt und zieht weg. Wie alt sind die anderen vier im Mittel? Steigt oder sinkt das Mittel?\karo{3}
\fusshilfe{\mbox{15:~$11$ Jahre, es sinkt}}
\end{pfaufg}

% L4-D: Mittel rückwärts und Änderungen
\begin{pfaufg}{}{16.}{}
Lea hat in vier Tests $14$, $17$, $12$ und $15$ Punkte. Im fünften Test gibt es höchstens $20$ Punkte. Am Ende will sie im Durchschnitt $15$ Punkte haben. Kann sie das noch schaffen? Wie viele Punkte braucht sie mindestens?\karo{3}
\fusshilfe{\mbox{16:~Ja, mit mindestens $17$ Punkten}}
\end{pfaufg}

% L4-T: Probetest
\begin{pfaufg}{}{17.}{}
An sechs Tagen wird mittags die Temperatur gemessen (in °C): $21$, $17$, $24$, $19$, $16$, $22$. Bestimme den Median.\karo{3}
\fusshilfe{\mbox{17:~$20$\,°C}}
\end{pfaufg}

% L4-T: Probetest
\begin{pfaufg}{}{18.}{}
Eine Bücherei zählt die Besucher einer Woche. Im Durchschnitt kamen $120$ Besucher am Tag. Der Wert für Donnerstag fehlt.\par\vspace{3pt} \adjustbox{max width=\linewidth}{\begin{varwidth}{50cm}\sachtabelle{lcccccc}{Tag & Mo & Di & Mi & Do & Fr & Sa}{Besucher & $95$ & $130$ & $110$ & \leerzelle & $140$ & $150$}\end{varwidth}}\par\vspace{3pt} Berechne die fehlende Zahl und trage sie ein.\karo{3}
\fusshilfe{\mbox{18:~Do: $95$ Besucher}}
\end{pfaufg}

% L4-T: Probetest
\begin{pfaufg}{}{19.}{}
Neun Kinder nennen ihre Schuhgröße: $38$, $40$, $37$, $40$, $41$, $39$, $40$, $38$, $42$.\par Aussage 1: „Die Spannweite beträgt $5$.“\par Aussage 2: „Der Median beträgt $41$.“\par Eine Aussage ist falsch. Kreuze sie an und berichtige sie.\par\smallskip \gkreuz{Aussage 1}\gkreuz{Aussage 2}
\fusshilfe{\mbox{19:~Aussage 2; richtig: Median $40$}}
\end{pfaufg}

% L4-T: Probetest
\begin{pfaufg}{}{20.}{}
Ein Chor hat $11$ Mitglieder. Ihr Alter in Jahren: $47$, $39$, $61$, $52$, $68$, $35$, $58$, $44$, $24$, $29$, $49$. a)~Berechne das Durchschnittsalter. b)~Die älteste und die jüngste Person verlassen den Chor. Ändert sich dadurch das Durchschnittsalter? Begründe.\karo{3}
\fusshilfe{\mbox{20:~\mbox{a) $46$ Jahre} \mbox{b) Nein, bleibt $46$ Jahre}}}
\end{pfaufg}

\par\vfill\noindent{\footnotesize Vorher: Streifen- und Kreisdiagramm\quad\textbullet\quad Weiter: Diagramme beurteilen und Boxplot\par}
\end{document}
